\section*{Step 48: Fork Selected Over Ladder}

Assume the Step47 common-carrier recognition premise: QM and GR are readouts of
a common carrier \(L\), warranted semiclassically by one \(\psi\) sourcing both
Born and stress-energy descents.

The ladder test asks whether \(q_{\rm GR}\) is a lawful quotient of
\(q_{\rm QM}\). On the finite four-mode carrier,
\[
q_{\rm QM}=(d_0,d_1,d_2), \qquad q_{\rm GR}=(d_0,d_2,d_3).
\]
The best linear quotient has positive residual, and the finite fiber test has
eight defect pairs. Adding the weak-RT/QM-accessible shadow
\(A_{\rm RT}=d_0+2d_2\) does not change the fiber obstruction. This is the
WEAK-RT ladder only: \(A_{\rm RT}\) is a function of QM-accessible modes and is
therefore structurally unable to recover \(d_3\). It does not test strong bulk
reconstruction from richer quantum data; that stronger claim is a named
residual, not refuted here.

The fork closes because both \(q_{\rm QM}\) and \(q_{\rm GR}\) descend from
\(L=(d_0,d_1,d_2,d_3)\) with zero residual. A nested control
\(q_{\rm GR}^{\rm nested}=(d_0,d_2,d_1)\) flips the verdict: the ladder closes
there with zero defect. Therefore the fork selection tracks the non-nested
mode structure rather than an always-on obstruction.

The Batch-A RT result is consequently read as a relation between the two
co-sourced arms of the fork, not as a quotient derivation of the full GR arm
from the QM arm.
